\section{Experimental Evaluation}
\label{sec:experiments}
\draftpending

% Target length: ~8 pages. BLOCKED until trained weights are available.
% Strategy: write captions and narrative structure NOW, drop in numbers and
% figures once the final checkpoint is evaluated.
% Tooling already in place: eval_checkpoints.py, testing/evaluate_policy.py
% (--axis fault and other OFAT axes, --seeds, --auto-nav).

\subsection{Experimental setup}
\label{subsec:exp-setup}
\draftpending
% TODO: checkpoint under test (best.pt of the final curriculum); fixed seed
% sets; episodes per configuration (~200-500); hardware; metrics: success
% rate, time-to-find (steps over successful episodes), coverage %, redundant
% visit ratio / overlap, per-drone contribution; confidence intervals via
% seed repetition.

\subsection{Baselines}
\label{subsec:exp-baselines}
\draftpending
% TODO: random policy; BFS auto-nav (deterministic sweep reference — NOT an
% upper bound, discuss why); optionally single drone vs team (value of
% cooperation).

\subsection{Generalization across randomization axes}
\label{subsec:exp-generalization}
\draftpending
% TODO: OFAT sweeps via testing/evaluate_policy.py — vision radius, comm
% range, team size, obstacle density; one curve per axis; in-range vs
% out-of-range probes.

\subsection{Fault robustness}
\label{subsec:exp-faults}
\draftpending
% TODO (headline experiment): success/time vs fault_prob (including p=0
% reference); behavioural adaptation: trajectories before/after a teammate's
% crash; does a survivor re-cover the dead drone's sector?; effect of comm
% range on how fast fault knowledge spreads (belief lag vs performance).

\subsection{Ablations}
\label{subsec:exp-ablations}
\draftpending
% TODO (time permitting): no communication (comm range 0); context without
% n_alive_belief; shared vs individual terminal reward.

\subsection{Qualitative analysis}
\label{subsec:exp-qualitative}
\draftpending
% TODO: trajectory plots, visit heatmaps, dispersion snapshots; link back to
% Sec. 6.6 emergent signatures (coverage-overlap measurements).

\subsection{Discussion and limitations}
\label{subsec:exp-discussion}
\draftpending
% TODO: simulation-only; grid abstraction; i.i.d. fault model; no adversary;
% sim-to-real gap.

% TODO: run experiments once weights are ready, then draft.
